Un ciclotró és un accelerador de partícules format per dos elèctrodes buits semicirculars (en forma de D) on actua un camp magnètic homogeni $\vec{B}$ perpendicular al pla horitzontal (pla de la figura). Així, a l'interior dels elèctrodes les partícules carregades positives, que es mouen en el pla horitzontal, descriuen una trajectòria circular. A l'espai buit que separa els dos elèctrodes s'aplica un camp elèctric altern $\vec{E}$, de manera que les partícules són accelerades. Inicialment, les partícules tenen poca velocitat i a cada cicle, en passar d'un semicercle a l'altre, van augmentant de velocitat i de radi de gir fins que finalment surten fora del ciclotró.

\figura[0.4]{fig1.png}

\begin{apartat}{1,25}
Les partícules tenen una càrrega elèctrica positiva $q$ i una massa $m$. Deduïu l'expressió de la velocitat de les partícules en funció del quocient càrrega-massa ($q/m$), del radi $r$ de la trajectòria de les partícules i del mòdul del camp magnètic. Comproveu que el temps de recorregut dins una D no depèn de la velocitat de les partícules. Per què el camp elèctric ha de ser altern? Trobeu l'expressió de la freqüència del camp elèctric.
\end{apartat}

\begin{apartat}{1,25}
El ciclotró té un radi de 0,50 m i un camp magnètic de 0,20 T. Quan hi accelerem protons, quina velocitat tenen quan surten del ciclotró? Quina és la longitud d'ona associada a aquests protons? Quin radi mínim hauria de tenir el ciclotró per a considerar que els protons tenen velocitats relativistes (és a dir, un 10\% de la velocitat de la llum)?
\end{apartat}

\dades{%
  $m_\mathrm{p} = 1,67\times 10^{-27}\ \mathrm{kg}$.\\
  $|e| = 1,602\times 10^{-19}\ \mathrm{C}$.\\
  $c = 3,00\times 10^{8}\ \mathrm{m\,s^{-1}}$.\\
  $h = 6,63\times 10^{-34}\ \mathrm{J\,s}$.}
